\documentclass[12pt]{article}

\usepackage{amsmath}
\usepackage{amssymb}
\usepackage{enumitem}
\usepackage{mathtools}

\title{Dynamic Programming Problems}
\author{Algorithms Assignment}
\date{\today}

\begin{document}

\maketitle

\section*{Dynamic Programming Problems}

\begin{enumerate}[label=\textbf{6.\arabic*.}]
\setcounter{enumi}{6}

\item A subsequence is palindromic if it is the same whether read left to right or right to left. For
instance, the sequence
$A, C, G, T, G, T, C, A, A, A, A, T, C, G$
has many palindromic subsequences, including $A, C, G, C, A$ and $A, A, A, A$ (on the other hand,
the subsequence $A, C, T$ is not palindromic). Devise an algorithm that takes a sequence $x[1 \ldots n]$
and returns the (length of the) longest palindromic subsequence. Its running time should be
$O(n^2)$.

\setcounter{enumi}{16}

\item Given an unlimited supply of coins of denominations $x_1, x_2, \ldots, x_n$, we wish to make change for
a value $v$; that is, we wish to find a set of coins whose total value is $v$. This might not be possible:
for instance, if the denominations are 5 and 10 then we can make change for 15 but not for 12.
Give an $O(nv)$ dynamic-programming algorithm for the following problem.

\textbf{Input:} $x_1, \ldots, x_n; v$.

\textbf{Question:} Is it possible to make change for $v$ using coins of denominations $x_1, \ldots, x_n$?

\item Consider the following variation on the change-making problem (Exercise 6.17): you are given
denominations $x_1, x_2, \ldots, x_n$, and you want to make change for a value $v$, but you are allowed to
use each denomination at most once. For instance, if the denominations are $1, 5, 10, 20$, then you
can make change for $16 = 1 + 15$ and for $31 = 1 + 10 + 20$ but not for $40$ (because you can't use 20
twice).

\begin{figure}[h]
\centering
\begin{minipage}{0.45\textwidth}
\centering
\textbf{First Binary Search Tree:}
\begin{verbatim}
       end
      /    \
    do      then
   /  \     /   \
begin else  if   while
\end{verbatim}
\end{minipage}
\hfill
\begin{minipage}{0.45\textwidth}
\centering
\textbf{Second Binary Search Tree:}
\begin{verbatim}
       do
      /  \
   begin  while
          /  
         if   
        /  \
     else   then
            \
             end
\end{verbatim}
\end{minipage}
\caption{Two binary search trees for the keywords of a programming language.}
\label{fig:bst}
\end{figure}

\textbf{Input:} Positive integers $x_1, x_2, \ldots, x_n$; another integer $v$.

\textbf{Output:} Can you make change for $v$, using each denomination $x_i$ at most once?

Show how to solve this problem in time $O(nv)$.

\item Here is yet another variation on the change-making problem (Exercise 6.17).
Given an unlimited supply of coins of denominations $x_1, x_2, \ldots, x_n$, we wish to make change for
a value $v$ using at most $k$ coins; that is, we wish to find a set of $\leq k$ coins whose total value is $v$.
This might not be possible: for instance, if the denominations are 5 and 10 and $k = 6$, then we
can make change for 55 but not for 65. Give an efficient dynamic-programming algorithm for the
following problem.

\textbf{Input:} $x_1, \ldots, x_n; k; v$.

\textbf{Question:} Is it possible to make change for $v$ using at most $k$ coins, of denominations
$x_1, \ldots, x_n$?

\item \textbf{Optimal binary search trees.} Suppose we know the frequency with which keywords occur in
programs of a certain language, for instance:
\begin{center}
\begin{tabular}{ll}
\texttt{begin} & 5\% \\
\texttt{do} & 40\% \\
\texttt{else} & 8\% \\
\texttt{end} & 4\% \\
\texttt{if} & 10\% \\
\texttt{then} & 10\% \\
\texttt{while} & 23\% \\
\end{tabular}
\end{center}

We want to organize them in a binary search tree, so that the keyword in the root is alphabetically
bigger than all the keywords in the left subtree and smaller than all the keywords in the right
subtree (and this holds for all nodes).

Figure 6.12 has a nicely-balanced example on the left. In this case, when a keyword is being
looked up, the number of comparisons needed is at most three: for instance, in finding \texttt{while},
only the three nodes \texttt{end}, \texttt{then}, and \texttt{while} get examined. But since we know the frequency
with which keywords are accessed, we can use an even more fine-tuned cost function, the average
number of comparisons to look up a word. For the search tree on the left, it is
\begin{align*}
\text{cost} &= 1(0.04) + 2(0.40 + 0.10) + 3(0.05 + 0.08 + 0.10 + 0.23)\\
&= 2.42
\end{align*}

By this measure, the best search tree is the one on the right, which has a cost of 2.18.

Give an efficient algorithm for the following task.

\textbf{Input:} $n$ words (in sorted order); frequencies of these words: $p_1, p_2, \ldots, p_n$.

\textbf{Output:} The binary search tree of lowest cost (defined above as the expected number
of comparisons in looking up a word).

\end{enumerate}

\end{document}
